%%%%%%%%%%%%%%%%%%%%%%%%%%%%%%%%%%%%%%%%%%%%%%%%%%%%%%%%%%%%%%%%%%%%%%%%
%                                                                      %
%     File: Thesis_Resumo.tex                                          %
%     Tex Master: Thesis.tex                                           %
%                                                                      %
%     Author: Andre C. Marta                                           %
%     Last modified :  4 Mar 2024                                      %
%                                                                      %
%%%%%%%%%%%%%%%%%%%%%%%%%%%%%%%%%%%%%%%%%%%%%%%%%%%%%%%%%%%%%%%%%%%%%%%%

\section*{Resumo}

% Add entry in the table of contents as section
\addcontentsline{toc}{section}{Resumo}

\emph{Descodificação sensível à qualidade} (QAD) — originalmente proposta para a tradução automática neuronal (MT) — é adaptada ao reconhecimento automático de fala (ASR) e à legendagem de imagens. Em vez de selecionar a hipótese de máxima verosimilhança através de descodificação/\emph{greedy beam search}, são gerados múltiplos candidatos e reordenados (\emph{reranked}) com recurso a métricas apropriadas à tarefa. É apresentada uma metodologia prática, juntamente com conjuntos de dados e um protocolo de avaliação, e é relacionado o comportamento empírico com leis de reordenação/\emph{reranking}.
\vfill

\textbf{\Large Palavras-chave:} Quality-Aware Decoding, Automatic Speech Recognition, Image Captioning, Minimum Bayes Risk

